%% kapa: exam=1 round=26 year=2015 period=2 form=A subject=accounting subject_ko=회계학 seq=13 number=53
%% source: Q-Net 감정평가사 기출문제 (cst00309) artlSeq=5138001 file=extracted/1008/5138001/제26회_감정평가사 1차_기출/제26회_감정평가사 1차_기출/2교시_감평_A형.pdf
%% conversion: kordoc 4.24.51 -> tools/gen1.py (auto); review=proofread-sonnet; flags=none
\begin{question}{53}
정부보조금의 회계처리에 관한 설명으로 옳지 않은 것은?
\begin{choices}
  \item 정부보조금은 정부보조금에 부수되는 조건의 준수와 보조금 수취에 대한 합리적인 확신이 있을 경우에만 정부보조금을 인식한다.
  \item 정부보조금으로 보전하려 하는 관련원가를 비용으로 인식하는 기간에 걸쳐 체계적인 기준에 따라 정부보조금을 당기손익으로 인식한다.
  \item 이미 발생한 비용이나 손실에 대한 보전 또는 향후의 관련원가 없이 기업에 제공되는 즉각적인 금융지원으로 수취하는 정부보조금은 정부보조금을 수취할 권리가 발생하는 기간에 당기손익으로 인식한다.
  \item 정부보조금은 비화폐성자산을 기업이 사용하도록 이전하는 형식을 취할 수 있다. 이러한 상황에서는 비화폐성자산을 공정가치로 회계처리할 수 없다.
  \item 자산관련 정부보조금은 재무상태표에 이연수익으로 표시하거나 자산의 장부금액을 결정할 때 차감하여 표시한다.
\end{choices}
\end{question}
\answer{4}
